\chapter{第 1 章分层答案}
\label{app:ch01-solutions}

\section*{口述概念}
\begin{enumerate}[leftmargin=*]
  \item 定义规定本书如何使用一个符号或术语，本身不是待判真的经验陈述；命题在给定定义和假设后具有真假值。“$R_p=\sum_iw_ir_i$”选定了本章的组合简单收益口径。凸组合界则声称该对象满足上下界，必须由证明建立一般性，或被一个合法反例推翻。
  \item $A\Rightarrow B$ 表示 $A$ 对 $B$ 充分，$B$ 对 $A$ 必要。若“没有未来信息”是可信回测的必要条件，那么发现未来信息足以否决回测；但没有发现未来信息并不能证明成本、选择偏差与统计不确定性都合格。通过必要条件只意味着“尚未因此被否决”。
  \item 合法反例必须属于原定义域、满足原命题全部假设，却让结论失败。研究某一假设的必要性时，最好只删除这一条。$w=(2,0,0)$ 虽非标准凸权重并可能越界，但其权重和为 2，同时破坏归一化；它不能把越界归因于非负性，因为三个权重仍非负。
\end{enumerate}

\section*{笔试推导}
\begin{enumerate}[leftmargin=*]
  \item 否定式为
  \[
    \exists x\in\mathcal D:\quad A(x)\land\neg B(x).
  \]
  因此反例一方面必须满足 $A(x)$，另一方面必须让 $B(x)$ 为假。只给一个 $B$ 为假的对象却不检查 $A$，不能反驳原命题。
  \item 设 $m=\min_i r_i,M=\max_i r_i$。利用 $\sum_iw_i=1$，
  \[
    R_p-m=\sum_iw_i(r_i-m)\geq0,
    \qquad
    M-R_p=\sum_iw_i(M-r_i)\geq0.
  \]
  每项非负来自 $w_i\geq0$、$r_i-m\geq0$ 与 $M-r_i\geq0$。$R_p=m$ 当且仅当每个 $w_i>0$ 的资产都满足 $r_i=m$；上界取等同理要求所有正权重资产都取得 $M$。
  \item 用逆否命题。若存在 $w_j<0$，因为结论声称对所有收益向量成立，可以专门取 $r_j=0$、其余 $r_i=1$。此时 $m=0,M=1$，而
  \[
    R_p=\sum_{i\ne j}w_i=1-w_j>1=M,
  \]
  所以上界失败。“对所有”正是在允许选择这条见证时使用的；若结论只针对一个固定样本，就不能任意替换 $r$。
\end{enumerate}

\section*{数值编程}
\begin{enumerate}[leftmargin=*]
  \item 贡献为 $(0.010,-0.003,0.006)$，总和 $0.013$；改动后为 $0.035$，高于最大分量 $0.03$，因为有负权重。
  \item 在 $w=1/2$ 两线相交且最坏收益为 0。偏离交点时其中一条严格为负，因此最大值为 0。
  \item 和的方差为 $8+2\gamma$，对应标准差为 $\sqrt{10},\sqrt8,\sqrt6$；只有零协方差时可用独立方差相加。
\end{enumerate}

\section*{研究判断}
\begin{enumerate}[leftmargin=*]
  \item 可写成“可信回测 $\Rightarrow$ 已合理建模手续费”，即手续费条件是可信性的必要条件之一，而非充分条件。即使通过，仍可能存在未来信息、幸存者偏差、参数选择泄漏、容量不足或不恰当的统计检验。
  \item $100$ 美元是货币，$2\%$ 日收益是无量纲但带一日持有期，$15\%$ 年化波动率带年化时间尺度，三者不能直接相加。替代报告可分栏展示：美元 P\&L；以明确资本基数换算的同期收益；由明确日频样本、年化规则与假设得到的波动率；最后另报风险调整指标及其定义，不制造无单位依据的“综合分”。
  \item 一百万次随机搜索只涉及有限多个收益向量，无法覆盖全称命题。一般证明应使用权重非负与权重和为 1，将组合收益夹在各分量的最小值与最大值之间。删除非负性但保留归一化，取 $w'=(1.5,-0.5,0)$，本例得到 $0.035>0.03$。这个反例说明非负权重条件在普遍结论中的作用。
\end{enumerate}

\section*{基础衔接：独立计算}
展开得到 $4+4+2\times1=10$，标准差 $\sqrt{10}$。独立缩放得到方差 $8$、标准差 $\sqrt8$，遗漏两份协方差；所有数值须使用同一收益单位。
